\documentclass[12pt]{article}
\usepackage[UTF8]{inputenc}
\usepackage{xeCJK}
\setCJKmainfont{Sarabun}
\usepackage{enumitem}
\usepackage{geometry}
\geometry{a4paper, margin=2.5cm}
\usepackage{setspace}
\onehalfspacing

\begin{document}

\begin{center}
{\Large\textbf{กัมมันตภาพรังสี}}\\[6pt]
{\normalsize แบบฝึกหัดเพิ่มเติม}\\[6pt]
{\normalsize วิชาฟิสิกส์ ม.ปลาย บทที่ 27}\\[4pt]
{\footnotesize ทั้งหมด 30 ข้อ}
\end{center}
\hrule
\bigskip

\section*{แบบฝึกหัดเพิ่มเติม}
\begin{enumerate}[label=\arabic*., itemsep=0.5\baselineskip, topsep=0.5\baselineskip, leftmargin=*]
\item เลขมวล (mass number) A ของนิวเคลียร์มีความหมายอย่างไร?
\begin{enumerate}[label=\alph*)., itemsep=0.3\baselineskip, topsep=0.3\baselineskip]
\item ก\ 
\begin{minipage}[t]{0.9\linewidth}
\raggedright จำนวนโปรตอน
\end{minipage}
\item ข\ 
\begin{minipage}[t]{0.9\linewidth}
\raggedright จำนวนนิวตรอน
\end{minipage}
\item ค\ 
\begin{minipage}[t]{0.9\linewidth}
\raggedright จำนวนโปรตอน + นิวตรอน
\end{minipage}
\item ง\ 
\begin{minipage}[t]{0.9\linewidth}
\raggedright จำนวนอิเล็กตรออน
\end{minipage}
\end{enumerate}
\vspace{-0.3\baselineskip}
\begin{small}
\textit{คำตอบ: ค}
\end{small}

\item ²³⁸U และ ²³⁵U เป็นไอโซโทปของธาตุใด?
\begin{enumerate}[label=\alph*)., itemsep=0.3\baselineskip, topsep=0.3\baselineskip]
\item ก\ 
\begin{minipage}[t]{0.9\linewidth}
\raggedright ทอเรียม
\end{minipage}
\item ข\ 
\begin{minipage}[t]{0.9\linewidth}
\raggedright เรเดียม
\end{minipage}
\item ค\ 
\begin{minipage}[t]{0.9\linewidth}
\raggedright ยูเรเนียม
\end{minipage}
\item ง\ 
\begin{minipage}[t]{0.9\linewidth}
\raggedright พลูโทเนียม
\end{minipage}
\end{enumerate}
\vspace{-0.3\baselineskip}
\begin{small}
\textit{คำตอบ: ค}
\end{small}

\item รังสีแอลฟา (α) ประกอบด้วยอนุภาคอะไร?
\begin{enumerate}[label=\alph*)., itemsep=0.3\baselineskip, topsep=0.3\baselineskip]
\item ก\ 
\begin{minipage}[t]{0.9\linewidth}
\raggedright 2 โปรตอน + 2 นิวตรอน
\end{minipage}
\item ข\ 
\begin{minipage}[t]{0.9\linewidth}
\raggedright 2 โปรตอน + 1 นิวตรอน
\end{minipage}
\item ค\ 
\begin{minipage}[t]{0.9\linewidth}
\raggedright 1 โปรตอน + 2 นิวตรอน
\end{minipage}
\item ง\ 
\begin{minipage}[t]{0.9\linewidth}
\raggedright 1 โปรตอน + 1 นิวตรอน
\end{minipage}
\end{enumerate}
\vspace{-0.3\baselineskip}
\begin{small}
\textit{คำตอบ: ก}
\end{small}

\item สารกัมมันตรังสี A มีครึ่งชีวิต 8 วัน หลัง 32 วันจะเหลือกี่เปอร์เซ็นต์?
\begin{enumerate}[label=\alph*)., itemsep=0.3\baselineskip, topsep=0.3\baselineskip]
\item ก\ 
\begin{minipage}[t]{0.9\linewidth}
\raggedright 3.125\%
\end{minipage}
\item ข\ 
\begin{minipage}[t]{0.9\linewidth}
\raggedright 6.25\%
\end{minipage}
\item ค\ 
\begin{minipage}[t]{0.9\linewidth}
\raggedright 12.5\%
\end{minipage}
\item ง\ 
\begin{minipage}[t]{0.9\linewidth}
\raggedright 1.5625\%
\end{minipage}
\end{enumerate}
\vspace{-0.3\baselineskip}
\begin{small}
\textit{คำตอบ: ก}
\end{small}

\item รังสีเบต้าบวก (β⁺) ประกอบด้วยอนุภาคอะไร?
\begin{enumerate}[label=\alph*)., itemsep=0.3\baselineskip, topsep=0.3\baselineskip]
\item ก\ 
\begin{minipage}[t]{0.9\linewidth}
\raggedright อิเล็กตรออน
\end{minipage}
\item ข\ 
\begin{minipage}[t]{0.9\linewidth}
\raggedright โพซิตรอน
\end{minipage}
\item ค\ 
\begin{minipage}[t]{0.9\linewidth}
\raggedright โปรตอน
\end{minipage}
\item ง\ 
\begin{minipage}[t]{0.9\linewidth}
\raggedright นิวตรอน
\end{minipage}
\end{enumerate}
\vspace{-0.3\baselineskip}
\begin{small}
\textit{คำตอบ: ข}
\end{small}

\item 1 Curie (Ci) เท่ากับกี่ Becquerel (Bq)?
\begin{enumerate}[label=\alph*)., itemsep=0.3\baselineskip, topsep=0.3\baselineskip]
\item ก\ 
\begin{minipage}[t]{0.9\linewidth}
\raggedright 3.7 × 10⁷ Bq
\end{minipage}
\item ข\ 
\begin{minipage}[t]{0.9\linewidth}
\raggedright 3.7 × 10¹⁰ Bq
\end{minipage}
\item ค\ 
\begin{minipage}[t]{0.9\linewidth}
\raggedright 3.7 × 10⁹ Bq
\end{minipage}
\item ง\ 
\begin{minipage}[t]{0.9\linewidth}
\raggedright 3.7 × 10⁸ Bq
\end{minipage}
\end{enumerate}
\vspace{-0.3\baselineskip}
\begin{small}
\textit{คำตอบ: ข}
\end{small}

\item ไอโซโทปกัมมันตรังสี (radioisotope) ต่างจากไอโซโทปเสถียรอย่างไร?
\begin{enumerate}[label=\alph*)., itemsep=0.3\baselineskip, topsep=0.3\baselineskip]
\item ก\ 
\begin{minipage}[t]{0.9\linewidth}
\raggedright มีจำนวนโปรตอนมากกว่า
\end{minipage}
\item ข\ 
\begin{minipage}[t]{0.9\linewidth}
\raggedright มีนิวเคลียสที่ไม่เสถียรและสลายตัวเอง
\end{minipage}
\item ค\ 
\begin{minipage}[t]{0.9\linewidth}
\raggedright มีอิเล็กตรออนมากกว่า
\end{minipage}
\item ง\ 
\begin{minipage}[t]{0.9\linewidth}
\raggedright มีขนาดเล็กกว่า
\end{minipage}
\end{enumerate}
\vspace{-0.3\baselineskip}
\begin{small}
\textit{คำตอบ: ข}
\end{small}

\item ทำไม C-14 จึงมีอยู่ในชั้นบรรยากาศตลอดเวลา?
\begin{enumerate}[label=\alph*)., itemsep=0.3\baselineskip, topsep=0.3\baselineskip]
\item ก\ 
\begin{minipage}[t]{0.9\linewidth}
\raggedright เพราะมันเป็นธาตุเสถียร
\end{minipage}
\item ข\ 
\begin{minipage}[t]{0.9\linewidth}
\raggedright เพราะรังสีคอสมิกชนอะตอม N-14 แล้วเปลี่ยนเป็น C-14
\end{minipage}
\item ค\ 
\begin{minipage}[t]{0.9\linewidth}
\raggedright เพราะพืชดูด C-14 จากดิน
\end{minipage}
\item ง\ 
\begin{minipage}[t]{0.9\linewidth}
\raggedright เพราะ C-14 ถูกสร้างจากปฏิกิริยาเคมี
\end{minipage}
\end{enumerate}
\vspace{-0.3\baselineskip}
\begin{small}
\textit{คำตอบ: ข}
\end{small}

\item พลังงานยึดเหนี่ยวนิวเคลียร์ต่อนิวคลีออนสูงสุดพบที่ธาตุใด?
\begin{enumerate}[label=\alph*)., itemsep=0.3\baselineskip, topsep=0.3\baselineskip]
\item ก\ 
\begin{minipage}[t]{0.9\linewidth}
\raggedright ไฮโดรเจน
\end{minipage}
\item ข\ 
\begin{minipage}[t]{0.9\linewidth}
\raggedright ฮีเลียม
\end{minipage}
\item ค\ 
\begin{minipage}[t]{0.9\linewidth}
\raggedright เหล็ก
\end{minipage}
\item ง\ 
\begin{minipage}[t]{0.9\linewidth}
\raggedright ยูเรเนียม
\end{minipage}
\end{enumerate}
\vspace{-0.3\baselineskip}
\begin{small}
\textit{คำตอบ: ค}
\end{small}

\item ข้อใดไม่ใช่รังสีที่เกิดจากการสลายตัวของนิวเคลียร์กัมมันตรังสี?
\begin{enumerate}[label=\alph*)., itemsep=0.3\baselineskip, topsep=0.3\baselineskip]
\item ก\ 
\begin{minipage}[t]{0.9\linewidth}
\raggedright แอลฟา (α)
\end{minipage}
\item ข\ 
\begin{minipage}[t]{0.9\linewidth}
\raggedright เบต้า (β)
\end{minipage}
\item ค\ 
\begin{minipage}[t]{0.9\linewidth}
\raggedright แกมมา (γ)
\end{minipage}
\item ง\ 
\begin{minipage}[t]{0.9\linewidth}
\raggedright ไมโครเวฟ
\end{minipage}
\end{enumerate}
\vspace{-0.3\baselineskip}
\begin{small}
\textit{คำตอบ: ง}
\end{small}

\item สารกัมมันตรังสีมีค่าคงที่การสลายตัว λ = 0.05 วัน⁻¹ ครึ่งชีวิตเท่ากับเท่าใด?
\begin{enumerate}[label=\alph*)., itemsep=0.3\baselineskip, topsep=0.3\baselineskip]
\item ก\ 
\begin{minipage}[t]{0.9\linewidth}
\raggedright 5 วัน
\end{minipage}
\item ข\ 
\begin{minipage}[t]{0.9\linewidth}
\raggedright 13.86 วัน
\end{minipage}
\item ค\ 
\begin{minipage}[t]{0.9\linewidth}
\raggedright 14 วัน
\end{minipage}
\item ง\ 
\begin{minipage}[t]{0.9\linewidth}
\raggedright 20 วัน
\end{minipage}
\end{enumerate}
\vspace{-0.3\baselineskip}
\begin{small}
\textit{คำตอบ: ข}
\end{small}

\item สารกัมมันตรังสี 6 × 10²³ อะตอม มีครึ่งชีวิต 5 วินาที หลัง 20 วินาทีจะเหลือกี่อะตอม?
\begin{enumerate}[label=\alph*)., itemsep=0.3\baselineskip, topsep=0.3\baselineskip]
\item ก\ 
\begin{minipage}[t]{0.9\linewidth}
\raggedright 3 × 10²³
\end{minipage}
\item ข\ 
\begin{minipage}[t]{0.9\linewidth}
\raggedright 7.5 × 10²²
\end{minipage}
\item ค\ 
\begin{minipage}[t]{0.9\linewidth}
\raggedright 1.5 × 10²²
\end{minipage}
\item ง\ 
\begin{minipage}[t]{0.9\linewidth}
\raggedright 3.75 × 10²²
\end{minipage}
\end{enumerate}
\vspace{-0.3\baselineskip}
\begin{small}
\textit{คำตอบ: ง}
\end{small}

\item ธาตุกัมมันตรังสี Ra-226 สลายตัวเป็น Rn-222 โดยการปล่อยอนุภาคใด?
\begin{enumerate}[label=\alph*)., itemsep=0.3\baselineskip, topsep=0.3\baselineskip]
\item ก\ 
\begin{minipage}[t]{0.9\linewidth}
\raggedright β
\end{minipage}
\item ข\ 
\begin{minipage}[t]{0.9\linewidth}
\raggedright α
\end{minipage}
\item ค\ 
\begin{minipage}[t]{0.9\linewidth}
\raggedright γ
\end{minipage}
\item ง\ 
\begin{minipage}[t]{0.9\linewidth}
\raggedright นิวตรอน
\end{minipage}
\end{enumerate}
\vspace{-0.3\baselineskip}
\begin{small}
\textit{คำตอบ: ข}
\end{small}

\item ไอโซโทปไฮโดรเจนที่มี 1 โปรตอน 1 นิวตรอน เรียกว่าอะไร?
\begin{enumerate}[label=\alph*)., itemsep=0.3\baselineskip, topsep=0.3\baselineskip]
\item ก\ 
\begin{minipage}[t]{0.9\linewidth}
\raggedright ไฮโดรเจนธรรมดา (H-1)
\end{minipage}
\item ข\ 
\begin{minipage}[t]{0.9\linewidth}
\raggedright ดิวเทอเรียม (D หรือ H-2)
\end{minipage}
\item ค\ 
\begin{minipage}[t]{0.9\linewidth}
\raggedright ไตรเทียม (T หรือ H-3)
\end{minipage}
\item ง\ 
\begin{minipage}[t]{0.9\linewidth}
\raggedright ไฮโดรเจน-4
\end{minipage}
\end{enumerate}
\vspace{-0.3\baselineskip}
\begin{small}
\textit{คำตอบ: ข}
\end{small}

\item ถ้าซากดึกดำบรรพ์มี C-14 เหลือ 12.5\% ของค่าปกติ ซากดึกดำบรรพ์นี้มีอายุเท่าใด?
\begin{enumerate}[label=\alph*)., itemsep=0.3\baselineskip, topsep=0.3\baselineskip]
\item ก\ 
\begin{minipage}[t]{0.9\linewidth}
\raggedright ประมาณ 11,460 ปี
\end{minipage}
\item ข\ 
\begin{minipage}[t]{0.9\linewidth}
\raggedright ประมาณ 17,190 ปี
\end{minipage}
\item ค\ 
\begin{minipage}[t]{0.9\linewidth}
\raggedright ประมาณ 22,920 ปี
\end{minipage}
\item ง\ 
\begin{minipage}[t]{0.9\linewidth}
\raggedright ประมาณ 5,730 ปี
\end{minipage}
\end{enumerate}
\vspace{-0.3\baselineskip}
\begin{small}
\textit{คำตอบ: ข}
\end{small}

\item กัมมันตภาพรังสี A ของสารกัมมันตรังสี 1 Ci มีค่าเท่าใดในหน่วย Bq?
\begin{enumerate}[label=\alph*)., itemsep=0.3\baselineskip, topsep=0.3\baselineskip]
\item ก\ 
\begin{minipage}[t]{0.9\linewidth}
\raggedright 3.7 × 10⁷ Bq
\end{minipage}
\item ข\ 
\begin{minipage}[t]{0.9\linewidth}
\raggedright 3.7 × 10⁹ Bq
\end{minipage}
\item ค\ 
\begin{minipage}[t]{0.9\linewidth}
\raggedright 3.7 × 10¹⁰ Bq
\end{minipage}
\item ง\ 
\begin{minipage}[t]{0.9\linewidth}
\raggedright 3.7 × 10⁸ Bq
\end{minipage}
\end{enumerate}
\vspace{-0.3\baselineskip}
\begin{small}
\textit{คำตอบ: ค}
\end{small}

\item ทำไมรังสีแอลฟาจึงมีอำนาจทะลุทะลวงต่ำที่สุดในบรรดารังสีทั้ง 3 ชนิด?
\begin{enumerate}[label=\alph*)., itemsep=0.3\baselineskip, topsep=0.3\baselineskip]
\item ก\ 
\begin{minipage}[t]{0.9\linewidth}
\raggedright เพราะมันหนักมากและเคลื่อนที่ช้า
\end{minipage}
\item ข\ 
\begin{minipage}[t]{0.9\linewidth}
\raggedright เพราะมันมีประจุไฟฟ้าสูง ทำให้ถูกดูดกลืนง่าย
\end{minipage}
\item ค\ 
\begin{minipage}[t]{0.9\linewidth}
\raggedright เพราะมันเป็นคลื่น
\end{minipage}
\item ง\ 
\begin{minipage}[t]{0.9\linewidth}
\raggedright เพราะมันมีขนาดเล็กที่สุด
\end{minipage}
\end{enumerate}
\vspace{-0.3\baselineskip}
\begin{small}
\textit{คำตอบ: ข}
\end{small}

\item พลังงานยึดเหนี่ยวนิวเคลียร์ของ ²³⁵U มีค่าประมาณเท่าใดต่อนิวคลีออน?
\begin{enumerate}[label=\alph*)., itemsep=0.3\baselineskip, topsep=0.3\baselineskip]
\item ก\ 
\begin{minipage}[t]{0.9\linewidth}
\raggedright 7.6 MeV
\end{minipage}
\item ข\ 
\begin{minipage}[t]{0.9\linewidth}
\raggedright 8.8 MeV
\end{minipage}
\item ค\ 
\begin{minipage}[t]{0.9\linewidth}
\raggedright 9.5 MeV
\end{minipage}
\item ง\ 
\begin{minipage}[t]{0.9\linewidth}
\raggedright 6.5 MeV
\end{minipage}
\end{enumerate}
\vspace{-0.3\baselineskip}
\begin{small}
\textit{คำตอบ: ก}
\end{small}

\item เมื่อ Co-60 (Z=27) สลายตัวโดยปล่อย β⁻ ได้นิกเกิล-60 หรือไม่?
\begin{enumerate}[label=\alph*)., itemsep=0.3\baselineskip, topsep=0.3\baselineskip]
\item ก\ 
\begin{minipage}[t]{0.9\linewidth}
\raggedright ได้ เพราะ Z เพิ่ม 1 จาก 27 เป็น 28 ซึ่งคือ Ni
\end{minipage}
\item ข\ 
\begin{minipage}[t]{0.9\linewidth}
\raggedright ไม่ได้ เพราะ A ไม่คงที่
\end{minipage}
\item ค\ 
\begin{minipage}[t]{0.9\linewidth}
\raggedright ได้ เพราะ N เพิ่ม 1
\end{minipage}
\item ง\ 
\begin{minipage}[t]{0.9\linewidth}
\raggedright ไม่ได้ เพราะ Ni-60 ไม่มีอยู่จริง
\end{minipage}
\end{enumerate}
\vspace{-0.3\baselineskip}
\begin{small}
\textit{คำตอบ: ก}
\end{small}

\item สารกัมมันตรังสีมีกัมมันตภาพรังสีเริ่มต้น A₀ = 800 Bq หลัง 30 วัน กัมมันตภาพรังสีเหลือ 100 Bq ครึ่งชีวิตเท่ากับเท่าใด?
\begin{enumerate}[label=\alph*)., itemsep=0.3\baselineskip, topsep=0.3\baselineskip]
\item ก\ 
\begin{minipage}[t]{0.9\linewidth}
\raggedright 10 วัน
\end{minipage}
\item ข\ 
\begin{minipage}[t]{0.9\linewidth}
\raggedright 15 วัน
\end{minipage}
\item ค\ 
\begin{minipage}[t]{0.9\linewidth}
\raggedright 20 วัน
\end{minipage}
\item ง\ 
\begin{minipage}[t]{0.9\linewidth}
\raggedright 7.5 วัน
\end{minipage}
\end{enumerate}
\vspace{-0.3\baselineskip}
\begin{small}
\textit{คำตอบ: ก}
\end{small}

\item หลักฐานทางโบราณคดีพบว่าซากดึกดำบรรพ์มีอัตราส่วน C-14/C-12 ลดลง 75\% จากค่าปกติ ซากดึกดำบรรพ์มีอายุเท่าใด?
\begin{enumerate}[label=\alph*)., itemsep=0.3\baselineskip, topsep=0.3\baselineskip]
\item ก\ 
\begin{minipage}[t]{0.9\linewidth}
\raggedright ประมาณ 2,300 ปี
\end{minipage}
\item ข\ 
\begin{minipage}[t]{0.9\linewidth}
\raggedright ประมาณ 5,700 ปี
\end{minipage}
\item ค\ 
\begin{minipage}[t]{0.9\linewidth}
\raggedright ประมาณ 11,400 ปี
\end{minipage}
\item ง\ 
\begin{minipage}[t]{0.9\linewidth}
\raggedright ประมาณ 1,400 ปี
\end{minipage}
\end{enumerate}
\vspace{-0.3\baselineskip}
\begin{small}
\textit{คำตอบ: ค}
\end{small}

\item สารกัมมันตรังสีมีจำนวนนิวเคลียสเริ่มต้น N₀ = 10¹⁰ และมีครึ่งชีวิต 8 วัน หลัง 24 วันจะเหลือนิวเคลียสกี่อะตอม?
\begin{enumerate}[label=\alph*)., itemsep=0.3\baselineskip, topsep=0.3\baselineskip]
\item ก\ 
\begin{minipage}[t]{0.9\linewidth}
\raggedright 6.25 × 10⁸
\end{minipage}
\item ข\ 
\begin{minipage}[t]{0.9\linewidth}
\raggedright 6.25 × 10⁷
\end{minipage}
\item ค\ 
\begin{minipage}[t]{0.9\linewidth}
\raggedright 1.25 × 10⁹
\end{minipage}
\item ง\ 
\begin{minipage}[t]{0.9\linewidth}
\raggedright 2.5 × 10⁸
\end{minipage}
\end{enumerate}
\vspace{-0.3\baselineskip}
\begin{small}
\textit{คำตอบ: ก}
\end{small}

\item นิวเคลียร์ K-40 (Z=19, A=40) สลายตัวทั้งแบบ β⁻ และ β⁺ (electron capture) ผลิตภัณฑ์ที่ได้คืออะไร?
\begin{enumerate}[label=\alph*)., itemsep=0.3\baselineskip, topsep=0.3\baselineskip]
\item ก\ 
\begin{minipage}[t]{0.9\linewidth}
\raggedright Ca-40 และ Ar-40
\end{minipage}
\item ข\ 
\begin{minipage}[t]{0.9\linewidth}
\raggedright Ca-40 เท่านั้น
\end{minipage}
\item ค\ 
\begin{minipage}[t]{0.9\linewidth}
\raggedright Ar-40 เท่านั้น
\end{minipage}
\item ง\ 
\begin{minipage}[t]{0.9\linewidth}
\raggedright Sc-40 และ Ca-40
\end{minipage}
\end{enumerate}
\vspace{-0.3\baselineskip}
\begin{small}
\textit{คำตอบ: ก}
\end{small}

\item สารกัมมันตรังสีมี λ = 0.2 วินาที⁻¹ หลัง 5 วินาที จะเหลือกี่เปอร์เซ็นต์ของจำนวนเริ่มต้น?
\begin{enumerate}[label=\alph*)., itemsep=0.3\baselineskip, topsep=0.3\baselineskip]
\item ก\ 
\begin{minipage}[t]{0.9\linewidth}
\raggedright e⁻¹ ≈ 37\%
\end{minipage}
\item ข\ 
\begin{minipage}[t]{0.9\linewidth}
\raggedright e⁻² ≈ 14\%
\end{minipage}
\item ค\ 
\begin{minipage}[t]{0.9\linewidth}
\raggedright e⁻⁰.⁵ ≈ 61\%
\end{minipage}
\item ง\ 
\begin{minipage}[t]{0.9\linewidth}
\raggedright e⁻¹.⁵ ≈ 22\%
\end{minipage}
\end{enumerate}
\vspace{-0.3\baselineskip}
\begin{small}
\textit{คำตอบ: ง}
\end{small}

\item ข้อใดไม่ใช่คุณสมบัติของนิวตรอน?
\begin{enumerate}[label=\alph*)., itemsep=0.3\baselineskip, topsep=0.3\baselineskip]
\item ก\ 
\begin{minipage}[t]{0.9\linewidth}
\raggedright มีมวลใกล้เคียงโปรตอน
\end{minipage}
\item ข\ 
\begin{minipage}[t]{0.9\linewidth}
\raggedright ไม่มีประจุไฟฟ้า
\end{minipage}
\item ค\ 
\begin{minipage}[t]{0.9\linewidth}
\raggedright สลายตัวเองเป็นโปรตอน + อิเล็กตรออน + แอนตินิวตริโน
\end{minipage}
\item ง\ 
\begin{minipage}[t]{0.9\linewidth}
\raggedright มีประจุบวก 1
\end{minipage}
\end{enumerate}
\vspace{-0.3\baselineskip}
\begin{small}
\textit{คำตอบ: ง}
\end{small}

\item ไอโซโทป H-3 (ไตรเทียม) มีโปรตอน นิวตรอน และอิเล็กตรออนกี่ตัว?
\begin{enumerate}[label=\alph*)., itemsep=0.3\baselineskip, topsep=0.3\baselineskip]
\item ก\ 
\begin{minipage}[t]{0.9\linewidth}
\raggedright 1p, 2n, 1e
\end{minipage}
\item ข\ 
\begin{minipage}[t]{0.9\linewidth}
\raggedright 1p, 3n, 1e
\end{minipage}
\item ค\ 
\begin{minipage}[t]{0.9\linewidth}
\raggedright 3p, 1n, 3e
\end{minipage}
\item ง\ 
\begin{minipage}[t]{0.9\linewidth}
\raggedright 1p, 1n, 1e
\end{minipage}
\end{enumerate}
\vspace{-0.3\baselineskip}
\begin{small}
\textit{คำตอบ: ก}
\end{small}

\item การสลายตัวแบบแกมมา (γ) มักเกิดร่วมกับการสลายตัวแบบ α หรือ β เพราะเหตุใด?
\begin{enumerate}[label=\alph*)., itemsep=0.3\baselineskip, topsep=0.3\baselineskip]
\item ก\ 
\begin{minipage}[t]{0.9\linewidth}
\raggedright เพราะรังสี γ เกิดจากการเปลี่ยนแปลงของนิวตรอน
\end{minipage}
\item ข\ 
\begin{minipage}[t]{0.9\linewidth}
\raggedright เพราะนิวเคลียร์หลังจากปล่อย α หรือ β มักอยู่ในสถานะกระตุ้นแล้วปล่อย γ
\end{minipage}
\item ค\ 
\begin{minipage}[t]{0.9\linewidth}
\raggedright เพราะรังสี γ ต้องมีพลังงานสูงมาก
\end{minipage}
\item ง\ 
\begin{minipage}[t]{0.9\linewidth}
\raggedright เพราะการสลาย γ เป็นการเปลี่ยนธาตุ
\end{minipage}
\end{enumerate}
\vspace{-0.3\baselineskip}
\begin{small}
\textit{คำตอบ: ข}
\end{small}

\item สารกัมมันตรังสีมีกัมมันตภาพรังสี 1 MBq หมายความว่าอย่างไร?
\begin{enumerate}[label=\alph*)., itemsep=0.3\baselineskip, topsep=0.3\baselineskip]
\item ก\ 
\begin{minipage}[t]{0.9\linewidth}
\raggedright มี 1 ล้านนิวเคลียส
\end{minipage}
\item ข\ 
\begin{minipage}[t]{0.9\linewidth}
\raggedright มีการสลายตัว 1 ล้านครั้งต่อวินาที
\end{minipage}
\item ค\ 
\begin{minipage}[t]{0.9\linewidth}
\raggedright มีครึ่งชีวิต 1 ล้านวินาที
\end{minipage}
\item ง\ 
\begin{minipage}[t]{0.9\linewidth}
\raggedright มีนิวเคลียส 1 ล้านกรัม
\end{minipage}
\end{enumerate}
\vspace{-0.3\baselineskip}
\begin{small}
\textit{คำตอบ: ข}
\end{small}

\item ข้อจำกัดของการหาอายุด้วยคาร์บอน-14 คืออะไร?
\begin{enumerate}[label=\alph*)., itemsep=0.3\baselineskip, topsep=0.3\baselineskip]
\item ก\ 
\begin{minipage}[t]{0.9\linewidth}
\raggedright ใช้ได้กับทุกวัตถุโบราณ
\end{minipage}
\item ข\ 
\begin{minipage}[t]{0.9\linewidth}
\raggedright ใช้ได้กับวัตถุที่มีอายุไม่เกิน \textasciitilde{}50,000 ปี และต้องมีคาร์บอนเป็นส่วนประกอบ
\end{minipage}
\item ค\ 
\begin{minipage}[t]{0.9\linewidth}
\raggedright ใช้ได้กับทุกชนิดของไอโซโทปคาร์บอน
\end{minipage}
\item ง\ 
\begin{minipage}[t]{0.9\linewidth}
\raggedright ไม่มีข้อจำกัด
\end{minipage}
\end{enumerate}
\vspace{-0.3\baselineskip}
\begin{small}
\textit{คำตอบ: ข}
\end{small}

\item ทำไมการสลายตัวของนิวเคลียร์จึงเป็นกระบวนการสุ่ม (random)?
\begin{enumerate}[label=\alph*)., itemsep=0.3\baselineskip, topsep=0.3\baselineskip]
\item ก\ 
\begin{minipage}[t]{0.9\linewidth}
\raggedright เพราะไม่ทราบกฎของธรรมชาติที่ควบคุม
\end{minipage}
\item ข\ 
\begin{minipage}[t]{0.9\linewidth}
\raggedright เพราะไม่สามารถทำนายได้ว่านิวเคลียสตัวใดจะสลายตัวก่อน แม้ว่าทุกตัวมีโอกาสเท่ากัน
\end{minipage}
\item ค\ 
\begin{minipage}[t]{0.9\linewidth}
\raggedright เพราะนิวเคลียสทุกตัวมีอายุเท่ากัน
\end{minipage}
\item ง\ 
\begin{minipage}[t]{0.9\linewidth}
\raggedright เพราะรังสีมีผลต่อการสลายตัว
\end{minipage}
\end{enumerate}
\vspace{-0.3\baselineskip}
\begin{small}
\textit{คำตอบ: ข}
\end{small}

\end{enumerate}
\end{document}